\documentclass[11pt,a4paper]{article}

% ----------------------------------
% Packages
% ----------------------------------
\usepackage[utf8]{inputenc}
\usepackage[T1]{fontenc}
\usepackage[a4paper,top=2.1cm,bottom=2.1cm,left=2.25cm,right=2.25cm]{geometry}
\usepackage{xcolor}
\usepackage{hyperref}
\usepackage{enumitem}
\usepackage{booktabs}
\usepackage{array}
\usepackage{tabularx}
\usepackage{longtable}
\usepackage{titlesec}
\usepackage{fancyhdr}
\usepackage{listings}
\usepackage{listingsutf8}
\usepackage{lmodern}
\usepackage{microtype}
\usepackage{float}
\usepackage{upquote}

% ----------------------------------
% Document information
% ----------------------------------
\newcommand{\DocumentType}{Vulnerability Notes}
\newcommand{\VulnerabilityTitle}{API Testing}
\newcommand{\SourceName}{PortSwigger Web Security Academy}
\newcommand{\SourceURL}{https://portswigger.net/web-security/api-testing}
\newcommand{\AuthorName}{Youssef Reda}

% ----------------------------------
% Colors
% ----------------------------------
\definecolor{primary}{RGB}{25,55,95}
\definecolor{secondary}{RGB}{75,95,120}
\definecolor{linkblue}{RGB}{35,85,155}
\definecolor{codegray}{RGB}{95,95,95}
\definecolor{codebg}{RGB}{247,248,250}
\definecolor{codeframe}{RGB}{210,216,224}

% ----------------------------------
% Code listings
% ----------------------------------
\lstset{
    inputencoding=utf8,
    basicstyle=\ttfamily\small,
    backgroundcolor=\color{codebg},
    frame=single,
    rulecolor=\color{codeframe},
    framesep=7pt,
    xleftmargin=4pt,
    xrightmargin=4pt,
    aboveskip=8pt,
    belowskip=8pt,
    breaklines=true,
    breakatwhitespace=false,
    columns=fullflexible,
    keepspaces=true,
    showstringspaces=false,
    tabsize=2
}

\lstdefinelanguage{JSON}{
    basicstyle=\ttfamily\small,
    numbers=none,
    showstringspaces=false,
    morestring=[b]",
    stringstyle=\color{secondary}
}

\lstdefinelanguage{HTTP}{
    basicstyle=\ttfamily\small,
    numbers=none,
    showstringspaces=false,
    morestring=[b]",
    stringstyle=\color{secondary},
    morecomment=[l]{\#}
}

\lstdefinelanguage{JavaScript}{
    language=Java,
    basicstyle=\ttfamily\small,
    numbers=none,
    showstringspaces=false,
    morekeywords={const,let,var,async,await,fetch,return,import,export,from},
    keywordstyle=\color{primary}\bfseries,
    stringstyle=\color{secondary},
    morestring=[b]',
    morestring=[b]"
}

% ----------------------------------
% Hyperlinks and PDF metadata
% ----------------------------------
\hypersetup{
    colorlinks=true,
    linkcolor=primary,
    urlcolor=linkblue,
    citecolor=primary,
    pdftitle={\VulnerabilityTitle},
    pdfauthor={\AuthorName}
}

\urlstyle{same}

% ----------------------------------
% Paragraphs, lists, and headings
% ----------------------------------
\setlength{\parindent}{0pt}
\setlength{\parskip}{4pt}

\setlist[itemize]{
    leftmargin=1.5em,
    itemsep=2pt,
    topsep=3pt
}

\setlist[enumerate]{
    leftmargin=1.8em,
    itemsep=2pt,
    topsep=3pt
}

\titleformat{\section}
  {\Large\bfseries\color{primary}}{}{0pt}{}

\titleformat{\subsection}
  {\large\bfseries\color{secondary}}{}{0pt}{}

\titleformat{\subsubsection}
  {\normalsize\bfseries}{}{0pt}{}

\titlespacing*{\section}{0pt}{16pt}{6pt}
\titlespacing*{\subsection}{0pt}{11pt}{4pt}
\titlespacing*{\subsubsection}{0pt}{8pt}{3pt}

% ----------------------------------
% Unnumbered headings in Table of Contents
% ----------------------------------
\newcommand{\nsection}[1]{%
  \section*{#1}%
  \phantomsection%
  \addcontentsline{toc}{section}{#1}%
}

\newcommand{\nsubsection}[1]{%
  \subsection*{#1}%
  \phantomsection%
  \addcontentsline{toc}{subsection}{#1}%
}

\newcommand{\nsubsubsection}[1]{%
  \subsubsection*{#1}%
  \phantomsection%
  \addcontentsline{toc}{subsubsection}{#1}%
}

\setcounter{tocdepth}{2}

% ----------------------------------
% Header and footer
% ----------------------------------
\pagestyle{fancy}
\fancyhf{}

\fancyhead[L]{\small\textcolor{secondary}{Web Security Notes}}
\fancyhead[R]{\small\textit{API Testing}}

\fancyfoot[C]{\thepage}

\renewcommand{\headrulewidth}{0.4pt}
\renewcommand{\footrulewidth}{0pt}

\setlength{\headheight}{14pt}

% ----------------------------------
% Document
% ----------------------------------
\begin{document}

% ----------------------------------
% Cover page
% ----------------------------------
\begin{titlepage}
    \thispagestyle{empty}
    \centering

    \vspace*{1.8cm}
    {\Large\color{secondary}\bfseries \DocumentType\par}

    \vspace{1.2cm}
    {\Huge\bfseries\color{primary}\VulnerabilityTitle\par}

    \vspace{0.8cm}
    {\large Concepts, Attack Techniques, Detection, and Defenses\par}

    \vfill

    \begin{flushleft}
        \textbf{Reference}\\
        \href{\SourceURL}{\SourceName}

        \vspace{0.6cm}
        \textbf{Prepared by}\\
        \AuthorName
    \end{flushleft}
\end{titlepage}

% ----------------------------------
% Table of Contents
% ----------------------------------
\pagenumbering{roman}
\tableofcontents

\clearpage
\pagenumbering{arabic}

% ----------------------------------
% Notes
% ----------------------------------

\nsection{OWASP API Security Top 10 Mapping}

The OWASP API Security Top 10 identifies critical security risks specific to APIs. Many of these risks overlap with vulnerabilities covered in existing Web Security topics.
The following table maps each OWASP API risk to its relevant Web Security Academy topics.

\begin{longtable}{p{0.39\textwidth} p{0.55\textwidth}}
    \toprule
    \textbf{OWASP API Risk} & \textbf{Relevant Web Security Academy Topics} \\
    \midrule
    \endfirsthead

    \toprule
    \textbf{OWASP API Risk} & \textbf{Relevant Web Security Academy Topics} \\
    \midrule
    \endhead

    API1: Broken Object Level Authorization
                            &
    Access Control Vulnerabilities and Privilege Escalation
    \\

    API2: Broken Authentication
                            &
    Authentication Vulnerabilities; OAuth 2.0 Authentication Vulnerabilities; JWT Attacks
    \\

    API3: Broken Object Property Level Authorization
                            &
    Mass Assignment Vulnerabilities
    \\

    API4: Unrestricted Resource Consumption
                            &
    Race Conditions; File Upload Vulnerabilities
    \\

    API5: Broken Function Level Authorization
                            &
    Access Control Vulnerabilities and Privilege Escalation
    \\

    API6: Unrestricted Access to Sensitive Business Flows
                            &
    Business Logic Vulnerabilities
    \\

    API7: Server Side Request Forgery
                            &
    Server-Side Request Forgery (SSRF)
    \\

    API8: Security Misconfiguration
                            &
    Cross-Origin Resource Sharing (CORS); Information Disclosure; HTTP Host Header Attacks; HTTP Request Smuggling
    \\

    API9: Improper Inventory Management
                            &
    API Testing
    \\

    API10: Unsafe Consumption of APIs
                            &
    API Testing
    \\

    \bottomrule
\end{longtable}

\newpage
\nsection{API Recon}

API reconnaissance is the process of gathering information about an API to identify its attack surface and understand how to interact with its endpoints.

\nsubsection{Identifying API Endpoints}

An API endpoint is a URL path where an API receives requests to interact with a specific resource.

Example:
\begin{lstlisting}[language=HTTP]
GET /api/books HTTP/1.1
Host: example.com
\end{lstlisting}

The endpoint is \texttt{/api/books}, which retrieves a list of books.

During reconnaissance, identify the following information:
\begin{itemize}
    \item \textbf{Input Data:} Required and optional parameters processed by the API.
    \item \textbf{HTTP Methods:} Supported methods, such as GET, POST, PUT, PATCH, and DELETE.
    \item \textbf{Media Formats:} Supported formats, such as JSON and XML.
    \item \textbf{Rate Limits:} Restrictions on the number of requests.
    \item \textbf{Authentication:} Authentication mechanisms required to access endpoints.
\end{itemize}

\nsubsection{API Documentation}

API documentation describes how developers can interact with and integrate APIs.

There are two main types:
\begin{itemize}
    \item \textbf{Human-readable:} Designed for developers, providing explanations, examples, and usage scenarios.
    \item \textbf{Machine-readable:} Structured formats that can be processed automatically, such as JSON and XML.
\end{itemize}

OpenAPI is a common specification for describing APIs, typically using JSON or YAML.

Always review publicly available API documentation as an initial reconnaissance step.

\newpage
\nsubsection{Discovering API Documentation}

Even when documentation is not publicly advertised, it may be accessible through the application that uses the API.

Discovery techniques include:
\begin{itemize}
    \item Use Burp Scanner to crawl the application.
    \item Browse the application manually using Burp Browser and inspect HTTP history.
    \item Search for common documentation paths.
    \item Use Burp Intruder with a wordlist of common paths.
\end{itemize}

Common documentation paths:
\begin{lstlisting}
 /api
 /swagger/index.html
 /openapi.json
\end{lstlisting}

\textbf{Parent Path Enumeration:}

When an endpoint is discovered, investigate its parent paths to identify additional resources or documentation.

For example, if the discovered endpoint is:
\begin{lstlisting}
/api/swagger/v1/users/123
\end{lstlisting}

Also investigate:
\begin{lstlisting}
/api/swagger/v1
/api/swagger
/api
\end{lstlisting}

\nsubsection{Analyzing Machine-readable Documentation}

Machine-readable API documentation can be analyzed using automated tools to identify endpoints and understand their expected parameters and request formats.

Useful tools include:
\begin{itemize}
    \item \textbf{Burp Scanner:} Crawls and audits OpenAPI documentation and other JSON or YAML documentation.
    \item \textbf{OpenAPI Parser BApp:} Parses OpenAPI specifications within Burp Suite.
    \item \textbf{Postman:} Imports API definitions and helps construct and send HTTP requests.
    \item \textbf{SoapUI:} Tests APIs, particularly SOAP-based services.
\end{itemize}

\newpage
\nsection{Identifying API Endpoints}

\nsubsection{Discovering Endpoints Through the Application}

Browsing applications that use an API can reveal additional endpoints, even when API documentation is available, as documentation may be outdated or inaccurate.

Useful techniques include:
\begin{itemize}
    \item \textbf{Burp Scanner:} Crawl the application to discover endpoints.
    \item \textbf{Burp Browser:} Manually browse the application and inspect HTTP traffic.
    \item \textbf{JavaScript Analysis:} Review JavaScript files for API endpoint references that have not been triggered through normal browsing.
\end{itemize}

Burp Scanner automatically extracts some endpoints from JavaScript files. For more extensive extraction, use the JS Link Finder BApp or manually inspect the files.

\nsubsection{Interacting with API Endpoints}

Use Burp Repeater and Burp Intruder to interact with discovered endpoints and investigate their behavior.

Test different:
\begin{itemize}
    \item HTTP methods.
    \item Content types.
    \item Parameters and input values.
\end{itemize}

Carefully examine HTTP responses and error messages, as they may disclose useful information about expected request formats, parameters, and input validation.

\nsubsection{Identifying Supported HTTP Methods}

HTTP methods specify the actions performed on resources.

Common methods include:
\begin{itemize}
    \item \textbf{GET:} Retrieves a resource.
    \item \textbf{POST:} Submits data or creates a resource.
    \item \textbf{PUT:} Creates or replaces a resource.
    \item \textbf{PATCH:} Applies partial modifications.
    \item \textbf{DELETE:} Deletes a resource.
    \item \textbf{OPTIONS:} Retrieves information about supported communication options, including HTTP methods.
\end{itemize}

A single endpoint may support multiple methods, exposing different functionalities.

Example:
\begin{lstlisting}[language=HTTP]
GET /api/tasks
POST /api/tasks
DELETE /api/tasks/1
\end{lstlisting}

Use Burp Intruder's built-in HTTP verbs list to automatically test different methods against an endpoint.
\textbf{Note:} Test potentially destructive methods against low-priority objects to avoid unintended modifications, deletions, or excessive record creation.

\nsubsection{Identifying Supported Content Types}

APIs may process input differently depending on the request's Content-Type header.

Changing the Content-Type and reformatting the request body accordingly may help to:
\begin{itemize}
    \item Trigger errors that disclose useful information.
    \item Bypass flawed security controls.
    \item Identify differences in input processing and validation logic.
\end{itemize}

For example, an API may validate JSON input correctly but process XML using an insecure parser, potentially exposing vulnerabilities such as XXE.
The Content Type Converter BApp can convert request bodies between JSON and XML formats.

\nsubsection{Using Intruder to Find Hidden Endpoints}

Burp Intruder can be used to enumerate hidden endpoints by replacing known path components with common function names.

For example, if the following endpoint is discovered:
\begin{lstlisting}[language=HTTP]
PUT /api/user/update
\end{lstlisting}

Place \texttt{update} in an Intruder payload position and test alternative names:

\begin{lstlisting}
/api/user/delete
/api/user/add
/api/user/create
/api/user/view
\end{lstlisting}

Use wordlists based on:
\begin{itemize}
    \item Common API naming conventions.
    \item Industry-specific terminology.
    \item Application functionality discovered during initial reconnaissance.
\end{itemize}

Application-specific wordlists are generally more useful than relying entirely on generic endpoint lists.

\newpage
\nsection{Finding Hidden Parameters}

\nsubsection{Hidden Parameter Discovery}

APIs may support undocumented parameters that are not exposed through the frontend or API documentation. Discovering these parameters can reveal additional functionality and potential vulnerabilities.

Useful Burp Suite tools include:
\begin{itemize}
    \item \textbf{Burp Intruder:} Uses wordlists of common and application-specific parameter names to replace existing parameters or add new ones.
    \item \textbf{Param Miner BApp:} Automatically guesses up to 65,536 parameter names per request, using information from the target scope.
    \item \textbf{Content Discovery:} Discovers resources and content that are not linked from visible application content, including parameters.
\end{itemize}

\nsubsection{Mass Assignment Vulnerabilities}

Mass Assignment (also known as Auto-binding) occurs when a framework automatically binds request parameters to fields of an internal object.
If the application does not properly restrict which fields users can modify, attackers may modify unintended properties, potentially leading to privilege escalation or unauthorized data modification.

\nsubsubsection{Identifying Hidden Parameters}

Inspect API responses, particularly GET responses, to identify object properties that are not included in normal update requests.

Example update request:
\begin{lstlisting}[language=JSON]
{
    "username": "wiener",
    "email": "wiener@example.com"
}
\end{lstlisting}

Corresponding GET response:
\begin{lstlisting}[language=JSON]
{
    "id": 123,
    "name": "John Doe",
    "email": "john@example.com",
    "isAdmin": false
}
\end{lstlisting}

Properties such as \texttt{id} and \texttt{isAdmin} may be candidates for hidden parameter testing, as they appear in the returned object but are absent from the normal update request.

\newpage
\nsubsubsection{Testing Mass Assignment}

To test whether hidden properties are bound to internal objects:

\textbf{1. Submit the parameter with its current valid value:}

\begin{lstlisting}[language=JSON]
{
    "username": "wiener",
    "email": "wiener@example.com",
    "isAdmin": false
}
\end{lstlisting}

\textbf{2. Submit an invalid value:}

\begin{lstlisting}[language=JSON]
{
    "username": "wiener",
    "email": "wiener@example.com",
    "isAdmin": "foo"
}
\end{lstlisting}

Observe differences in responses and application behavior. Different responses may indicate that the parameter is being processed by the application.

\textbf{3. Attempt to modify the property:}

\begin{lstlisting}[language=JSON]
{
    "username": "wiener",
    "email": "wiener@example.com",
    "isAdmin": true
}
\end{lstlisting}

Verify the result by retrieving the user object again and checking whether the property was actually modified.
Finally, test whether the modified privileges provide access to administrative functionality.
A successful modification of restricted properties may result in privilege escalation.

\newpage
\nsection{Preventing Vulnerabilities in APIs}

\nsubsection{API Security Best Practices}

When designing APIs, security should be considered from the beginning. Apply the following practices to reduce the API attack surface and prevent common vulnerabilities.

\begin{itemize}
    \item \textbf{Secure API Documentation:} Restrict access to API documentation when the API is not intended to be publicly accessible.
    \item \textbf{Keep Documentation Updated:} Ensure documentation accurately reflects the API's current endpoints, HTTP methods, and parameters, providing legitimate testers with complete visibility of the attack surface.
    \item \textbf{Allowlist HTTP Methods:} Explicitly permit only the HTTP methods required by each endpoint and reject unsupported methods.
    \item \textbf{Validate Content Types:} Verify that incoming requests and outgoing responses use the expected content types.
    \item \textbf{Use Generic Error Messages:} Avoid exposing internal implementation details, database errors, file paths, or sensitive information. Store technical details in server logs instead.
    \item \textbf{Protect All API Versions:} Apply appropriate security controls to every active API version, including legacy versions, rather than protecting only the current production version.
\end{itemize}

\nsubsection{Preventing Mass Assignment}

Mass assignment occurs when a framework automatically binds request parameters to internal object properties without sufficient restrictions.

To prevent mass assignment vulnerabilities:
\begin{itemize}
    \item \textbf{Property Allowlist (Recommended):} Explicitly define the properties users are permitted to create or update. Reject or ignore all other properties.
    \item \textbf{Property Blocklist:} Explicitly define sensitive properties that users must not modify, such as \texttt{isAdmin}, \texttt{role}, and \texttt{accountBalance}.
\end{itemize}

An allowlist is generally safer because only explicitly permitted properties can be modified. Newly introduced sensitive properties are not automatically exposed to user modification.

\textbf{Example:}

If users are only allowed to update their username and email, the application should accept:
\begin{lstlisting}[language=JSON]
{
    "username": "wiener",
    "email": "wiener@example.com"
}
\end{lstlisting}

Properties such as \texttt{isAdmin} and \texttt{role} should remain controlled by server-side logic and must not be modifiable through user requests.

\end{document}
